% Vietnamese translation for OI Public Library
% Source-PDF: 其他 OTHERS/一份不太简短的LATEX介绍.pdf
\documentclass[11pt]{article}
\usepackage{oipl-vn}

\usepackage{array}
\usepackage{longtable}
\usepackage{tikz}
\usetikzlibrary{arrows.meta,positioning}
\setmonofont{Noto Sans Mono}[Scale=MatchLowercase]

\newcommand{\Lshort}{\textit{The Not So Short Introduction to \LaTeXe}}
\newcommand{\CTAN}{\textsc{ctan}}
\newcommand{\BibTeX}{Bib\TeX}
\newcommand{\MakeIndex}{MakeIndex}
\newcommand{\cmd}[1]{\texttt{\textbackslash #1}}
\newcommand{\env}[1]{\texttt{#1}}
\newcommand{\pkg}[1]{\textsf{#1}}
\newcommand{\cls}[1]{\textsf{#1}}
\newcommand{\file}[1]{\texttt{#1}}
\newcommand{\meta}[1]{\textit{\textlangle #1\textrangle}}
\newtheorem{theorem}{Định lý}
\newtheorem{definition}{Định nghĩa}

\title{Một Bản Giới Thiệu Không Quá Ngắn Về \LaTeXe}
\author{Tobias Oetiker\\
Hubert Partl, Irene Hyna và Elisabeth Schlegl\\[0.4em]
Bản dịch tiếng Việt cho OI Public Library}
\date{Dựa trên bản 3.20, ngày 9 tháng 8 năm 2001}

\begin{document}
\maketitle

\renewcommand{\figurename}{Hình}

\origtitle{Một bản giới thiệu không quá ngắn về \LaTeXe, hay học \LaTeXe\
trong 93 phút}
\sourcepath{Others / A not so short introduction to LaTeX}

\begin{abstract}
Tài liệu này giới thiệu \LaTeXe\ cho người mới bắt đầu nhưng muốn đi đủ xa để
tự soạn một văn bản nghiêm túc. Nội dung đi từ cấu trúc tối thiểu của một tập
tin nguồn, cách \LaTeX\ xử lý chữ, đoạn, bảng, danh sách và tài liệu lớn, đến
công thức toán học, tài liệu tham khảo, chỉ mục, gói mở rộng, định nghĩa lệnh
mới và tinh chỉnh bố cục. Các đoạn mã được giữ ở dạng ví dụ \LaTeX\ có thể đọc
trực tiếp; bản dịch ưu tiên diễn đạt rõ ràng bằng tiếng Việt thay vì tái tạo
chính xác từng trang của bản in.
\end{abstract}

\section*{Bản quyền và giấy phép}
\addcontentsline{toc}{section}{Bản quyền và giấy phép}

Tobias Oetiker giữ bản quyền năm 2000 và các bản phát hành của \Lshort. Tài
liệu gốc là tài liệu tự do: người đọc có thể phát hành lại hoặc sửa đổi theo
các điều khoản của GNU General Public License do Free Software Foundation công
bố, bản 2 hoặc bất kỳ bản kế tiếp nào theo lựa chọn của người phát hành.

Tài liệu được phát hành với hy vọng hữu ích, nhưng không kèm bảo đảm nào, kể
cả bảo đảm ngầm định về khả năng thương mại hoặc sự phù hợp với một mục đích
cụ thể. Bản gốc yêu cầu người nhận xem văn bản GNU General Public License để
biết đầy đủ chi tiết; nếu không nhận được bản sao giấy phép, có thể liên hệ
Free Software Foundation, 675 Mass Ave, Cambridge, MA 02139, USA.

\section*{Lời cảm ơn của bản nguồn}
\addcontentsline{toc}{section}{Lời cảm ơn của bản nguồn}

Nhiều tư liệu trong phần giới thiệu này bắt nguồn từ tài liệu \LaTeX~2.09 bằng
tiếng Đức do ba tác giả ở Vienna và Graz biên soạn:
\begin{itemize}
  \item Hubert Partl, Zentraler Informatikdienst der Universitat fur
        Bodenkultur Wien;
  \item Irene Hyna, Bundesministerium fur Wissenschaft und Forschung Wien;
  \item Elisabeth Schlegl, Graz.
\end{itemize}

Bản tiếng Đức cập nhật cho \LaTeXe\ do Jorg Knappen thực hiện có thể tìm trong
cây \CTAN\ tại \file{CTAN:/tex-archive/info/lshort/german}.

Trong quá trình chuẩn bị tài liệu, Tobias Oetiker đã tham khảo ý kiến cộng
đồng \texttt{comp.text.tex}. Các phản hồi, hiệu đính, gợi ý và phần bổ sung
từ cộng đồng đã giúp tài liệu có hình thức hiện tại. Những người được bản gốc
ghi nhận gồm Rosemary Bailey, Friedemann Brauer, Jan Busa, Markus Bruhwiler,
David Carlisle, Jose Carlos Santos, Mike Chapman, Christopher Chin, Carl
Cerecke, Chris McCormack, Wim van Dam, Jan Dittberner, Michael John Downes,
David Dureisseix, Elliot, David Frey, Robin Fairbairns, Jorg Fischer, Erik
Frisk, Frank, Kasper B. Graversen, Alexandre Guimond, Cyril Goutte, Greg
Gamble, Neil Hammond, Rasmus Borup Hansen, Joseph Hilferty, Bjorn Hvittfeldt,
Martien Hulsen, Werner Icking, Jakob, Eric Jacoboni, Alan Jeffrey, Byron
Jones, David Jones, Johannes-Maria Kaltenbach, Michael Koundouros, Andrzej
Kawalec, Alain Kessi, Christian Kern, Jorg Knappen, Kjetil Kjernsmo, Maik
Lehradt, Alexander Mai, Martin Maechler, Aleksandar S Milosevic, Claus Malten,
Kevin Van Maren, Lenimar Nunes de Andrade, Hubert Partl, John Refling, Mike
Ressler, Brian Ripley, Young U. Ryu, Bernd Rosenlecher, Chris Rowley,
Hanspeter Schmid, Craig Schlenter, Christopher Sawtell, Geoffrey Swindale,
Josef Tkadlec, Didier Verna, Fabian Wernli, Carl-Gustav Werner, David
Woodhouse, Chris York, Fritz Zaucker, Rick Zaccone và Mikhail Zotov.

\section*{Ghi nhận bản dịch tiếng Trung}
\addcontentsline{toc}{section}{Ghi nhận bản dịch tiếng Trung}

Bản tiếng Trung do nhóm người dùng CTEX thực hiện. Công việc dịch được điều
phối bởi quản trị viên diễn đàn CTEX ``Jingdian Wenti'' và kéo dài gần mười
tháng. Bản nguồn ghi nhận các thành viên sau:

\begin{center}
\begin{tabular}{>{\ttfamily}p{0.24\linewidth}p{0.22\linewidth}p{0.32\linewidth}}
\textnormal{ID diễn đàn CTEX} & \textnormal{Phần dịch} & \textnormal{Tập tin nguồn}\\
\hline
Jingdian Wenti & Lời nói đầu & overview.tex\\
Gaoyuan zhi lang & Chương 1 & things.tex\\
controlong & Chương 2 & typeset.tex\\
cxterm & Chương 3 & math.tex, lssym.tex\\
aloft & Chương 4 & spec.tex\\
ganzhi & Chương 5 & custom.tex
\end{tabular}
\end{center}

\tableofcontents

\section*{Lời nói đầu}
\addcontentsline{toc}{section}{Lời nói đầu}

\LaTeX\ là một hệ thống dàn trang rất thích hợp để tạo ra tài liệu khoa học và
tài liệu toán học có chất lượng in ấn cao. Nó cũng dùng tốt cho thư từ, báo cáo,
bài viết, sách và nhiều loại văn bản khác. \LaTeX\ dùng \TeX\ làm động cơ định
dạng.

Mục tiêu của tài liệu này không phải là hướng dẫn cài đặt một hệ thống
\LaTeX. Ở hầu hết các hệ điều hành phổ biến đều có bản phân phối \TeX; trong
môi trường đại học hoặc phòng máy, hệ thống thường đã được cài sẵn. Điều quan
trọng hơn đối với người soạn thảo là hiểu cách viết tập tin nguồn để \LaTeX\
có thể xử lý nó.

Nên đọc các chương theo thứ tự. Chương đầu tạo khung khái niệm về \TeX,
\LaTeX, tập tin nguồn và lớp tài liệu. Chương hai đi vào dàn trang văn bản.
Chương ba trình bày công thức toán. Chương bốn giới thiệu hình vẽ, tài liệu
tham khảo, chỉ mục và một số gói tiện ích. Chương năm nói về việc tùy biến:
lệnh mới, môi trường mới, phông chữ, khoảng cách, bố cục trang và hộp.

Khi cần tài liệu hoặc gói mở rộng, điểm xuất phát tốt nhất là \CTAN:
\url{https://www.ctan.org}. Trong tài liệu gốc, đường dẫn dạng
\file{CTAN:/tex-archive/...} chỉ vị trí tương ứng trong cây lưu trữ \CTAN.

\section{Kiến thức cơ bản}

\subsection{Tên gọi trong trò chơi}

\subsubsection{\TeX}

\TeX\ là một chương trình dàn trang do Donald E. Knuth tạo ra để xử lý văn bản
và công thức toán học. Knuth bắt đầu viết động cơ \TeX\ năm 1977, khi thiết bị
in số bắt đầu đi vào ngành xuất bản. Mục tiêu trực tiếp của ông là đảo chiều
sự suy giảm chất lượng in ấn mà ông thấy trong sách và bài báo toán học của
chính mình.

Hệ \TeX\ mà ta dùng ngày nay được công bố năm 1982 và được chỉnh sửa đáng kể
năm 1989 để hỗ trợ tốt hơn ký tự 8-bit và nhiều ngôn ngữ. \TeX\ nổi tiếng vì
độ ổn định cao, chạy được trên nhiều loại máy tính và có rất ít lỗi. Số phiên
bản của \TeX\ tiến dần tới \(\pi\); trong bản nguồn của tài liệu này, số phiên
bản được ghi là \(3.14159\).

Tên \TeX\ được phát âm gần như ``tech'', với âm cuối giống chữ cái Hy Lạp
\(\chi\), không phải như chữ ``x'' tiếng Anh. Trong văn bản ASCII, tên này
thường được viết là \texttt{TeX}.

\TeX\ không phải là một trình soạn thảo văn bản kiểu ``thấy gì được nấy''.
Người dùng viết một tập tin nguồn chứa văn bản và lệnh đánh dấu, sau đó chạy
\TeX\ để sinh kết quả dàn trang. Cách làm này đòi hỏi một bước biên dịch, nhưng
đổi lại nó cho phép tài liệu có cấu trúc rõ ràng, tái lập được và dễ quản lý
khi lớn dần.

\subsubsection{\LaTeX}

\LaTeX\ là một gói macro đặt trên \TeX, ban đầu do Leslie Lamport phát triển.
Mục đích của nó là cho phép tác giả dùng một thiết kế trang chuyên nghiệp đã
được định nghĩa sẵn để dàn và in tài liệu chất lượng cao. Nếu \TeX\ là động cơ,
\LaTeX\ là bộ điều khiển giúp tác giả nói ở mức cao hơn: đây là tiêu đề, đây
là mục, đây là định lý, đây là tài liệu tham khảo. Người viết không cần tự tính
mọi khoảng trắng, cỡ chữ và vị trí dòng.

Năm 1994, nhóm \LaTeX3 do Frank Mittelbach dẫn dắt cập nhật hệ macro \LaTeX,
đưa vào nhiều cải tiến được chờ đợi và thống nhất lại các bản vá rời rạc từng
tích lũy sau \LaTeX~2.09. Phiên bản mới được gọi là \LaTeXe\ để phân biệt với
bản cũ. Trong văn bản ASCII, có thể viết \LaTeXe\ là \texttt{LaTeX2e}; cách
đọc thường gặp là ``Lay-tech two e'' hoặc ``Lah-tech two e''.

\begin{figure}[ht]
\centering
\begin{tikzpicture}[
  node distance=0.7cm and 0.95cm,
  box/.style={draw, rounded corners=1.5pt, align=center, minimum height=7mm,
    inner xsep=5pt, font=\small},
  file/.style={box, font=\small\ttfamily},
  tool/.style={box, fill=gray!8},
  >=Latex
]
\node[tool] (editor) {trình soạn\\thảo};
\node[file, right=of editor] (tex) {.tex};
\node[tool, right=of tex] (latex) {\LaTeXe};
\node[file, right=of latex] (dvi) {.dvi};
\node[tool, right=of dvi] (driver) {driver};
\node[box, right=of driver] (output) {máy in\\màn hình};

\node[box, above=of latex] (packages) {lớp, gói\\AMS, plain};
\node[file, below=of latex] (log) {.log};
\node[tool, below=of dvi] (metafont) {MetaFont};
\node[file, below=of metafont] (mf) {.mf};
\node[file, right=of metafont] (tfm) {.tfm};
\node[file, right=of tfm] (pk) {.pk};

\draw[->] (editor) -- (tex);
\draw[->] (tex) -- (latex);
\draw[->] (latex) -- (dvi);
\draw[->] (dvi) -- (driver);
\draw[->] (driver) -- (output);
\draw[->] (packages) -- (latex);
\draw[->] (latex) -- (log);
\draw[->] (mf) -- (metafont);
\draw[->] (metafont) -- (tfm);
\draw[->] (metafont) -- (pk);
\draw[->] (tfm) -- (latex);
\draw[->] (pk) -- (driver);
\end{tikzpicture}
\caption{Các thành phần chính của một hệ \TeX/\LaTeX.}
\end{figure}

\subsection{Nền tảng của \LaTeX}

\subsubsection{Tác giả, người thiết kế sách và thợ sắp chữ}

Khi viết sách bằng một trình soạn thảo thông thường, tác giả thường phải vừa
viết nội dung vừa quyết định hình thức: dòng này cỡ chữ bao nhiêu, tiêu đề cách
đoạn sau mấy điểm, danh sách thụt vào bao nhiêu. Công việc đó vốn thuộc về
người thiết kế sách và thợ sắp chữ.

\LaTeX\ khuyến khích tách hai việc này. Tác giả đánh dấu ý nghĩa logic của
văn bản, chẳng hạn \cmd{section}, \cmd{emph}, \env{itemize}, \env{table}. Lớp
tài liệu và các gói sẽ quyết định cách các phần ấy xuất hiện trên trang. Nhờ
vậy một văn bản có thể đổi từ dạng bài báo sang báo cáo hoặc sách với ít thay
đổi hơn nhiều so với việc chỉnh tay từng tiêu đề.

Cách làm này khác đáng kể với các hệ ``thấy gì được nấy'' (WYSIWYG, viết tắt
của \emph{What You See Is What You Get}) như nhiều trình xử lý văn bản thông
dụng. Trong những hệ đó, tác giả vừa nhập nội dung vừa tương tác trực tiếp với
diện mạo cuối cùng trên màn hình. Với \LaTeX, thông thường không thấy ngay kết
quả cuối trong lúc gõ nguồn; sau khi biên dịch, ta xem trước bản dàn trang, sửa
nguồn nếu cần, rồi biên dịch lại trước khi in hoặc xuất PDF.

Ưu điểm của cách \LaTeX\ là văn bản phải nói rõ cấu trúc logic. Điều này giúp
tránh tình trạng tài liệu trông bắt mắt ở vài trang riêng lẻ nhưng thiếu cấu
trúc nhất quán khi đọc toàn bộ.

\subsubsection{Thiết kế bố cục}

Thiết kế bố cục tốt không chỉ là làm trang giấy ``đẹp mắt''. Nó liên quan tới
độ dài dòng, khoảng cách dòng, tỷ lệ lề, vị trí tiêu đề, đánh số, đầu trang,
chân trang và nhiều chi tiết nhỏ khác. Nếu người viết không có kinh nghiệm
trình bày sách, quyết định tùy hứng thường dẫn tới trang in khó đọc.

Bản nguồn nhấn mạnh rằng tài liệu được viết để đọc, không phải để treo trong
phòng triển lãm. Vì vậy tính dễ đọc và dễ hiểu quan trọng hơn ấn tượng thị
giác tức thời. Chẳng hạn, cỡ chữ và cách đánh số tiêu đề phải giúp người đọc
nhìn ra cấu trúc chương mục; chiều dài dòng phải đủ ngắn để mắt không mỏi,
nhưng đủ dài để trang giấy không bị vụn.

\LaTeX\ cung cấp bố cục mặc định có tính bảo thủ. Nó không luôn hợp với mọi
nhu cầu, nhưng thường là điểm khởi đầu tốt hơn việc tự vẽ trang từ đầu. Khi cần
thay đổi, hãy thay đổi có chủ đích và giữ cấu trúc logic của tài liệu.

\subsubsection{Ưu điểm và giới hạn}

\LaTeX\ có nhiều ưu điểm so với các trình xử lý văn bản thông thường:
\begin{itemize}
  \item cung cấp thiết kế dàn trang chuyên nghiệp, giúp tài liệu có dáng dấp
  của một ấn phẩm đã được sắp chữ;
  \item dàn công thức toán học thuận tiện và có chất lượng rất cao;
  \item tác giả chỉ cần học một số lệnh tương đối dễ hiểu để mô tả cấu trúc
  logic của tài liệu, thay vì vá víu từng chi tiết trang;
  \item các cấu trúc phức tạp như chú thích chân trang, chỉ mục, mục lục, danh
  sách hình bảng và tài liệu tham khảo được tạo tương đối dễ dàng;
  \item có nhiều gói mở rộng miễn phí cho những việc không nằm trong lõi
  \LaTeX, chẳng hạn chèn đồ họa PostScript/PDF, quản lý tài liệu tham khảo
  theo chuẩn xuất bản hoặc hỗ trợ ngôn ngữ;
  \item \LaTeX\ khuyến khích tác giả viết tài liệu có cấu trúc tốt, vì chính
  cấu trúc ấy là đầu vào để hệ thống chọn bố cục phù hợp;
  \item \TeX, động cơ định dạng của \LaTeX, là phần mềm tự do và có tính khả
  chuyển cao, nên có thể chạy trên rất nhiều nền tảng phần cứng và hệ điều hành.
\end{itemize}

\LaTeX\ cũng có giới hạn. Bản nguồn nói đùa rằng có người có thể liệt kê hàng
trăm nhược điểm; những điểm thực tế hơn gồm:
\begin{itemize}
  \item chỉnh sâu một bố cục hoàn toàn mới có thể khó và tốn thời gian, dù các
  bố cục có sẵn cho phép thay đổi nhiều tham số;
  \item \LaTeX\ không phù hợp với tài liệu phi cấu trúc hoặc cố ý lộn xộn;
  \item đường cong học tập ban đầu có thật: những việc rất trực quan trong
  trình soạn thảo tương tác có thể cần một lệnh, một môi trường hoặc một gói;
  \item nếu cố ép \LaTeX\ bằng nhiều chỉnh sửa cục bộ thay vì dùng cấu trúc
  logic, tài liệu có thể trở nên khó bảo trì.
\end{itemize}

\subsection{Tập tin nguồn \LaTeX}

Một tập tin nguồn \LaTeX\ là văn bản thường. Bạn có thể dùng bất kỳ trình soạn
thảo nào lưu được văn bản UTF-8. Trong tập tin ấy, nội dung và lệnh được trộn
lẫn. Lệnh bắt đầu bằng dấu gạch chéo ngược.

\subsubsection{Khoảng trắng}

\LaTeX\ xử lý một hay nhiều dấu cách liên tiếp như một dấu cách. Một lần xuống
dòng trong nguồn cũng tương đương một dấu cách. Dòng trống mới tạo ra đoạn văn
mới.

\begin{verbatim}
This is one line.     Several spaces
in the input are treated as one.

An empty line starts a new paragraph.
\end{verbatim}

Nếu thật sự cần khoảng trắng không được ngắt dòng, dùng dấu ngã:
\begin{verbatim}
See Figure~\ref{fig:sample} and Chapter~2.
\end{verbatim}

\subsubsection{Ký tự đặc biệt}

Một số ký tự có ý nghĩa đặc biệt trong \LaTeX:
\[
\#\quad \$\quad \%\quad \&\quad \_\quad \{\quad \}\quad \sim\quad
\hat{\ }\quad \backslash
\]

Để in chúng như chữ thường, phần lớn có thể được viết bằng cách đặt dấu gạch
chéo ngược phía trước:
\begin{verbatim}
\# \$ \% \& \_ \{ \}
\end{verbatim}
Riêng dấu gạch chéo ngược nên in bằng \cmd{textbackslash}, dấu mũ bằng
\cmd{textasciicircum}, và dấu ngã bằng \cmd{textasciitilde}.

\subsubsection{Lệnh \LaTeX}

Lệnh \LaTeX\ có hai dạng chính. Dạng thứ nhất là dấu gạch chéo ngược theo sau
bởi một chuỗi chữ cái, chẳng hạn \cmd{section}. Tên lệnh kết thúc khi gặp ký tự
không phải chữ cái. Dạng thứ hai là dấu gạch chéo ngược theo sau bởi đúng một
ký tự không phải chữ cái, chẳng hạn \verb|\\| để xuống dòng.

Tham số bắt buộc đặt trong ngoặc nhọn, tham số tùy chọn đặt trong ngoặc vuông:
\begin{verbatim}
\section{Title}
\documentclass[11pt,a4paper]{article}
\end{verbatim}

\subsubsection{Chú thích}

Dấu phần trăm bắt đầu một chú thích. Mọi ký tự từ dấu \verb|%| đến hết dòng sẽ
bị bỏ qua.

\begin{verbatim}
This is text. % This comment is ignored.
\end{verbatim}

Chú thích cũng có thể dùng để ngăn khoảng trắng do xuống dòng trong nguồn, nhất
là khi định nghĩa macro.

\subsection{Cấu trúc của tập tin nguồn}

Một tài liệu \LaTeX\ tối thiểu có dạng:
\begin{verbatim}
\documentclass{article}

\begin{document}
A very small document.
\end{document}
\end{verbatim}

Phần trước \verb|\begin{document}| gọi là lời mở đầu. Tại đó ta chọn lớp tài
liệu, nạp gói, định nghĩa lệnh và thiết lập các tùy chọn toàn cục. Phần giữa
\verb|\begin{document}| và \verb|\end{document}| là thân tài liệu. Mọi thứ sau
\verb|\end{document}| bị bỏ qua.

Một ví dụ thực tế hơn:
\begin{verbatim}
\documentclass[11pt,a4paper]{article}
\usepackage{amsmath}

\title{A Small Example}
\author{John Smith}
\date{\today}

\begin{document}
\maketitle
\tableofcontents

\section{Introduction}
Hello \LaTeX.

\end{document}
\end{verbatim}

\subsection{Một phiên làm việc dòng lệnh điển hình}

Quy trình cổ điển là:
\begin{enumerate}
  \item Soạn tập tin \file{document.tex}.
  \item Chạy \texttt{latex document.tex} hoặc với hệ thống hiện đại:
        \texttt{lualatex document.tex}.
  \item Với quy trình cổ điển, \LaTeX\ sinh tập tin \file{document.dvi}. Có thể
        xem nó bằng một trình xem DVI, chẳng hạn \texttt{xdvi document.dvi}.
  \item Nếu cần in hoặc đưa sang hệ khác, có thể đổi DVI sang PostScript bằng
        một lệnh kiểu:
\begin{verbatim}
dvips -Pcmz document.dvi -o document.ps
\end{verbatim}
        Với các hệ hiện đại, \texttt{pdflatex}, \texttt{xelatex} hoặc
        \texttt{lualatex} thường sinh PDF trực tiếp.
  \item Nếu tài liệu có mục lục, tham chiếu chéo hoặc thư mục, chạy lại một
        hoặc nhiều lần để các số trang và nhãn ổn định.
  \item Xem kết quả PDF.
\end{enumerate}

Khi có lỗi, \LaTeX\ thường dừng tại vị trí gần nơi nó không hiểu nguồn. Thông
báo lỗi đôi khi khó đọc, nhưng dòng đầu tiên thường cho biết nguyên nhân chính:
thiếu dấu ngoặc, sai tên lệnh, môi trường chưa đóng, hoặc dùng ký tự đặc biệt
không được thoát.

\subsection{Bố cục tài liệu}

\subsubsection{Lớp tài liệu}

Lớp tài liệu được chọn bằng \cmd{documentclass}. Các lớp chuẩn thường gặp:

\begin{center}
\begin{tabular}{>{\ttfamily}ll}
article & bài báo, báo cáo ngắn, tài liệu kỹ thuật;\\
report  & báo cáo dài, luận văn nhỏ, có chương;\\
book    & sách, có mặt trước và mặt sau;\\
letter  & thư;\\
slides  & bản trình chiếu kiểu cũ.
\end{tabular}
\end{center}

Một số tùy chọn phổ biến:

\begin{center}
\begin{tabular}{>{\ttfamily}ll}
10pt, 11pt, 12pt & cỡ chữ cơ sở;\\
a4paper, letterpaper & khổ giấy; ngoài ra còn có a5paper, b5paper, legalpaper;\\
fleqn & đặt công thức hiển thị lệch trái thay vì căn giữa;\\
leqno & đặt số phương trình ở bên trái;\\
titlepage, notitlepage & tiêu đề tài liệu có bắt đầu trang riêng hay không;\\
onecolumn, twocolumn & một hoặc hai cột;\\
oneside, twoside & in một mặt hoặc hai mặt;\\
openright, openany & chương mới bắt đầu ở trang phải hoặc trang bất kỳ.
\end{tabular}
\end{center}

Ví dụ:
\begin{verbatim}
\documentclass[11pt,a4paper,twoside]{article}
\end{verbatim}

\subsubsection{Gói mở rộng}

Gói mở rộng được nạp bằng \cmd{usepackage}. Chúng bổ sung lệnh, môi trường hoặc
thiết lập mới. Ví dụ:

\begin{center}
\begin{tabular}{>{\ttfamily}ll}
amsmath & công thức toán nâng cao;\\
amssymb & ký hiệu toán bổ sung;\\
graphicx & chèn và biến đổi hình ảnh;\\
fontspec & dùng phông OpenType với XeLaTeX hoặc LuaLaTeX;\\
babel, polyglossia & hỗ trợ quy ước ngôn ngữ;\\
hyperref & liên kết trong PDF;\\
fancyhdr & tùy chỉnh đầu trang và chân trang;\\
ifthen & lệnh điều kiện kiểu ``nếu... thì... ngược lại...'';\\
makeidx & hỗ trợ tạo chỉ mục;\\
syntonly & chỉ kiểm tra cú pháp, không sinh bản dàn trang;\\
inputenc & khai báo mã hóa đầu vào trong các hệ \LaTeX\ truyền thống.
\end{tabular}
\end{center}

Tham số tùy chọn của gói đặt trước tên gói:
\begin{verbatim}
\usepackage[unicode]{hyperref}
\end{verbatim}

\subsection{Các loại tập tin của \LaTeX}

\LaTeX\ dùng và tạo ra nhiều loại tập tin:

\begin{center}
\begin{tabular}{>{\ttfamily}ll}
.tex & tập tin nguồn \LaTeX\ hoặc \TeX;\\
.sty & gói mở rộng, nạp bằng \cmd{usepackage};\\
.cls & lớp tài liệu;\\
.dtx & tập tin \TeX\ có tài liệu hóa, thường dùng để phát hành gói;\\
.ins & tập tin cài đặt đi kèm \file{.dtx}, dùng để trích xuất mã gói;\\
.dvi & kết quả ``độc lập thiết bị'' của quy trình \texttt{latex} cổ điển;\\
.aux & thông tin phụ cho tham chiếu chéo;\\
.toc & dữ liệu mục lục;\\
.lof, .lot & danh sách hình và danh sách bảng;\\
.log & nhật ký biên dịch;\\
.bib & cơ sở dữ liệu tài liệu tham khảo của \BibTeX;\\
.idx & dữ liệu chỉ mục thô do \LaTeX\ ghi ra;\\
.ind & chỉ mục đã xử lý, được đưa vào lần biên dịch sau;\\
.ilg & nhật ký của \MakeIndex.
\end{tabular}
\end{center}

Không nên chỉnh tay các tập tin sinh ra như \file{.aux}, \file{.toc} hoặc
\file{.log}. Khi gặp lỗi lạ do dữ liệu phụ cũ, xóa các tập tin sinh ra rồi
biên dịch lại thường là cách xử lý đơn giản.

\subsubsection{Kiểu trang}

Kiểu trang điều khiển đầu trang và chân trang. Các kiểu chuẩn gồm:

\begin{center}
\begin{tabular}{>{\ttfamily}ll}
plain & số trang ở chân trang, thường dùng mặc định;\\
headings & đầu trang chứa tiêu đề chương hoặc mục;\\
empty & không có đầu trang, chân trang và số trang.
\end{tabular}
\end{center}

Chọn kiểu trang bằng:
\begin{verbatim}
\pagestyle{headings}
\thispagestyle{empty}
\end{verbatim}

\subsection{Tài liệu lớn}

Tài liệu dài nên chia thành nhiều tập tin. Lệnh \cmd{input} chèn nội dung của
một tập tin tại vị trí hiện tại:
\begin{verbatim}
\input{chapters/introduction}
\end{verbatim}

Lệnh \cmd{include} phù hợp hơn cho các phần lớn như chương, vì nó bắt đầu trang
mới và tạo tập tin phụ riêng:
\begin{verbatim}
\include{chapter1}
\include{chapter2}
\end{verbatim}

Khi chỉ muốn biên dịch một vài phần trong tài liệu lớn, dùng:
\begin{verbatim}
\includeonly{chapter2,chapter5}
\end{verbatim}
trong lời mở đầu. Các tham chiếu tới phần không biên dịch vẫn có thể giữ đúng
nếu tập tin phụ của chúng đã được tạo từ lần biên dịch trước.

Khi tài liệu lớn và chỉ muốn kiểm tra nhanh cú pháp, có thể dùng gói
\pkg{syntonly}. \LaTeX\ sẽ đọc nguồn và kiểm tra lệnh, nhưng không sinh tập tin
kết quả dàn trang:
\begin{verbatim}
\usepackage{syntonly}
\syntaxonly
\end{verbatim}

\section{Dàn trang văn bản}

\subsection{Cấu trúc của tài liệu và ngôn ngữ}

\LaTeX\ xem tài liệu như một cấu trúc phân cấp. Với lớp \cls{article}, các cấp
thường dùng là:
\begin{verbatim}
\part{...}
\section{...}
\subsection{...}
\subsubsection{...}
\paragraph{...}
\subparagraph{...}
\end{verbatim}

Với \cls{report} và \cls{book}, có thêm \cmd{chapter}. Không nên tự gõ số mục
vào tiêu đề; \LaTeX\ tự đánh số và cập nhật mục lục. Nếu muốn tiêu đề không
đánh số, dùng dạng có sao:
\begin{verbatim}
\section*{Lời cảm ơn}
\end{verbatim}

Văn bản tốt không chỉ là chuỗi câu. Nó có đoạn, mục, ví dụ, định lý, hình,
bảng, chú thích và tham chiếu. Hãy dùng lệnh thể hiện vai trò của từng phần.
Ví dụ, dùng \cmd{emph} để nhấn mạnh thay vì chọn thủ công chữ nghiêng; nếu cách
nhấn mạnh thay đổi, toàn bộ tài liệu vẫn nhất quán.

\subsection{Ngắt dòng và ngắt trang}

\subsubsection{Sắp đoạn}

\LaTeX\ tự ngắt dòng trong đoạn văn. Khi cần buộc xuống dòng mà không tạo đoạn
mới, dùng \verb|\\| hoặc \cmd{newline}. Khi cần bắt đầu trang mới, dùng
\cmd{newpage}. Các lệnh mạnh hơn như \cmd{clearpage} còn buộc in hết các hình
và bảng đang chờ.

Thông thường không nên ép ngắt dòng thủ công trong văn bản chạy. Nếu một đoạn
bị dàn xấu, nguyên nhân thường là từ quá dài, công thức trong dòng quá cồng
kềnh, hoặc dùng lệnh tạo hộp không thể ngắt.

\subsubsection{Ngắt từ}

\LaTeX\ có thuật toán ngắt từ tự động. Với các từ đặc biệt, có thể chỉ rõ điểm
ngắt bằng:
\begin{verbatim}
hy\-phen\-a\-tion
\end{verbatim}

Danh sách quy tắc ngắt từ toàn cục đặt trong lời mở đầu:
\begin{verbatim}
\hyphenation{FORTRAN Hy-phen-a-tion}
\end{verbatim}

Dấu \cmd{mbox} giữ một cụm không bị ngắt dòng:
\begin{verbatim}
\mbox{Mr.~Smith}
\end{verbatim}

\subsection{Học cách dàn các chuỗi ký tự}

Một số chuỗi cần được gõ theo quy ước của \LaTeX:

\begin{center}
\begin{tabular}{>{\ttfamily}ll}
\textbackslash today & ngày hiện tại;\\
\textbackslash LaTeX & logo \LaTeX;\\
\textbackslash TeX & logo \TeX;\\
\textbackslash dots & dấu ba chấm đúng khoảng cách;\\
\textbackslash ldots & dấu ba chấm thấp.
\end{tabular}
\end{center}

Ví dụ:
\begin{verbatim}
Today is \today. I am learning \LaTeX\ldots
\end{verbatim}

\subsection{Ký tự và ký hiệu đặc biệt}

\subsubsection{Dấu ngoặc kép}

Trong nguồn \LaTeX\ truyền thống, dấu ngoặc kép mở và đóng không gõ bằng cùng
một ký tự. Dấu mở dùng hai dấu huyền, dấu đóng dùng hai dấu nháy đơn:
\begin{verbatim}
``A quoted sentence.''
\end{verbatim}
Với LuaLaTeX và phông Unicode, có thể gõ trực tiếp dấu ngoặc tiếng Việt nếu
quy trình biên dịch hỗ trợ tốt.

\subsubsection{Gạch nối, gạch ngang và dấu trừ}

\LaTeX\ phân biệt:
\begin{center}
\begin{tabular}{>{\ttfamily}ll}
- & gạch nối trong từ ghép;\\
-- & gạch ngang ngắn;\\
--- & gạch ngang dài;\\
\$-\$ & dấu trừ trong toán học.
\end{tabular}
\end{center}

\subsubsection{Dấu ngã}

Dấu ngã trong nguồn tạo khoảng trắng không ngắt dòng. Để in ký hiệu ngã, dùng
\cmd{textasciitilde} trong văn bản hoặc \cmd{sim} trong toán học:
\[
  a \sim b.
\]

\subsubsection{Ký hiệu độ}

Ký hiệu độ có thể viết trong toán học:
\begin{verbatim}
$-30\,^{\circ}\mathrm{C}$
\end{verbatim}
cho kết quả \(-30\,^{\circ}\mathrm{C}\).

\subsubsection{Dấu ba chấm}

Không nên gõ ba dấu chấm liên tiếp. Dùng \cmd{ldots} trong văn bản:
\begin{verbatim}
Not like this ... but like this \ldots
\end{verbatim}

\subsubsection{Liên tự}

\TeX\ tự tạo một số liên tự như \textit{ff}, \textit{fi}, \textit{fl} tùy
phông chữ. Nếu cần ngăn liên tự tại ranh giới từ hoặc hình vị, dùng nhóm rỗng:
\begin{verbatim}
shel{}fful
\end{verbatim}

\subsubsection{Dấu phụ và ký tự đặc biệt}

Trong \LaTeX\ cổ điển, các chữ có dấu được tạo bằng lệnh như:
\begin{center}
\begin{tabular}{>{\ttfamily}ll}
\textbackslash '\{e\} & acute: \'{e};\\
\textbackslash `\{e\} & grave: \`{e};\\
\textbackslash \^{}\{o\} & circumflex: \^{o};\\
\textbackslash "\{u\} & umlaut: \"{u};\\
\textbackslash \~{}\{n\} & tilde: \~{n};\\
\textbackslash c\{c\} & cedilla: \c{c}.
\end{tabular}
\end{center}

Với LuaLaTeX, cách tự nhiên nhất cho tiếng Việt là lưu nguồn UTF-8 và gõ trực
tiếp chữ có dấu.

\subsection{Hỗ trợ ngôn ngữ quốc tế}

Một tài liệu đa ngôn ngữ cần nhiều hơn phông chữ. Mỗi ngôn ngữ có quy tắc ngắt
từ, tên tự động cho mục lục hoặc hình bảng, kiểu ngày tháng và dấu câu riêng.
Trong hệ \LaTeX\ hiện đại, \pkg{babel} hoặc \pkg{polyglossia} thường đảm nhiệm
việc này.

Ví dụ với \pkg{babel}:
\begin{verbatim}
\usepackage[vietnamese,english]{babel}
\end{verbatim}

Ví dụ với \pkg{polyglossia} trong LuaLaTeX:
\begin{verbatim}
\usepackage{polyglossia}
\setmainlanguage{vietnamese}
\setotherlanguage{english}
\end{verbatim}

\subsubsection{Hỗ trợ tiếng Đức}

Bản gốc dành một phần nhỏ cho tiếng Đức vì tài liệu đầu tiên được viết bằng
tiếng Đức. Tiếng Đức cần xử lý dấu umlaut, chữ \ss{} và quy tắc ngắt từ. Với
các hệ hiện đại, nạp ngôn ngữ Đức qua \pkg{babel} hoặc \pkg{polyglossia} sẽ
giải quyết phần lớn chi tiết này.

\subsection{Khoảng cách giữa các từ}

\TeX\ thường đặt khoảng cách sau dấu chấm câu cuối câu rộng hơn khoảng cách
giữa từ. Nó đoán dấu chấm cuối câu dựa trên chữ cái trước dấu chấm. Sau chữ
viết hoa, \TeX\ giả định đó là viết tắt nên không tăng khoảng cách:
\begin{verbatim}
Mr. Smith was happy.
I watched NASA\@. It launched a rocket.
\end{verbatim}

Dấu gạch chéo ngược trước khoảng trắng tạo khoảng trắng thường:
\begin{verbatim}
Prof.\ Smith
\end{verbatim}
Dấu \verb|~| tạo khoảng trắng không ngắt dòng:
\begin{verbatim}
Figure~3
\end{verbatim}

\subsection{Tiêu đề, chương và mục}

Thông tin tiêu đề đặt bằng \cmd{title}, \cmd{author} và \cmd{date}; lệnh
\cmd{maketitle} in ra tiêu đề.

\begin{verbatim}
\title{A Document}
\author{The Author}
\date{\today}

\begin{document}
\maketitle
\end{document}
\end{verbatim}

Mục lục được tạo bằng \cmd{tableofcontents}. Sau lần biên dịch đầu tiên, dữ
liệu mục lục mới được ghi vào tập tin phụ; vì vậy cần biên dịch lại để số trang
đúng.

\subsection{Tham chiếu chéo}

Đặt nhãn bằng \cmd{label} và tham chiếu bằng \cmd{ref} hoặc \cmd{pageref}.
\begin{verbatim}
\section{Algorithms}\label{sec:algo}
See Section~\ref{sec:algo} on page~\pageref{sec:algo}.
\end{verbatim}

Nhãn nên có tiền tố giúp nhận biết loại đối tượng, chẳng hạn \file{sec:},
\file{fig:}, \file{tab:}, \file{eq:}. Sau khi thêm hoặc đổi nhãn, cần biên
dịch lại đủ lần để tham chiếu cập nhật.

\subsection{Chú thích chân trang}

Chú thích chân trang viết bằng \cmd{footnote}:
\begin{verbatim}
This is a remark\footnote{This is a footnote.}.
\end{verbatim}

Không nên lạm dụng chú thích. Nếu thông tin đủ quan trọng, đặt nó vào thân văn
bản; nếu chỉ là chi tiết phụ, chú thích là nơi thích hợp.

\subsection{Nhấn mạnh}

Lệnh \cmd{emph} biểu diễn nhấn mạnh logic:
\begin{verbatim}
This is \emph{important}.
\end{verbatim}

Trong văn bản thường, kết quả thường là chữ nghiêng. Nếu dùng \cmd{emph} bên
trong đoạn đã nghiêng, nó có thể quay lại chữ đứng. Đây là lý do nên dùng
\cmd{emph} thay vì tự chọn \cmd{textit}.

\subsection{Môi trường}

Môi trường có dạng:
\begin{verbatim}
\begin{environment-name}
  ...
\end{environment-name}
\end{verbatim}

\subsubsection{\env{itemize}, \env{enumerate} và \env{description}}

\begin{verbatim}
\begin{itemize}
  \item An item.
  \item Another item.
\end{itemize}

\begin{enumerate}
  \item The first item.
  \item The second item.
\end{enumerate}

\begin{description}
  \item[TeX] A typesetting engine.
  \item[LaTeX] A higher-level macro package.
\end{description}
\end{verbatim}

\subsubsection{\env{flushleft}, \env{flushright} và \env{center}}

Các môi trường này căn trái, căn phải hoặc căn giữa một khối:
\begin{verbatim}
\begin{center}
Centered text.
\end{center}
\end{verbatim}

\subsubsection{\env{quote}, \env{quotation} và \env{verse}}

\env{quote} phù hợp cho trích dẫn ngắn. \env{quotation} dùng cho trích dẫn dài
có nhiều đoạn. \env{verse} dành cho thơ, nơi xuống dòng có ý nghĩa.

\subsubsection{In nguyên văn}

\env{verbatim} in nội dung gần như đúng như nguồn, rất hữu ích cho mã lệnh:
\begin{quote}
\ttfamily
\textbackslash begin\{verbatim\}\\
for i := 1 to 10 do\\
\mbox{\quad}writeln(i);\\
\textbackslash end\{verbatim\}
\end{quote}

Dạng trong dòng là \cmd{verb}:
\begin{verbatim}
\verb|printf("%d\n", x);|
\end{verbatim}

\subsubsection{Bảng}

Môi trường \env{tabular} tạo bảng:
\begin{verbatim}
\begin{tabular}{|r|l|}
\hline
7C0 & hexadecimal \\
3700 & octal \\
11111000000 & binary \\
\hline
1984 & decimal \\
\hline
\end{tabular}
\end{verbatim}

Các ký tự trong mô tả cột gồm \texttt{l} căn trái, \texttt{c} căn giữa,
\texttt{r} căn phải, \texttt{p\{}\meta{width}\texttt{\}} tạo cột đoạn văn có độ rộng
cố định, và \texttt{|} kẻ đường dọc. Dùng \cmd{hline} để kẻ đường ngang.

\subsection{Vật thể nổi}

Hình và bảng lớn không nên bị ép đứng đúng tại vị trí trong nguồn. \LaTeX\ dùng
vật thể nổi để đặt chúng ở nơi hợp lý hơn. Hai môi trường chuẩn là \env{figure}
và \env{table}.

\begin{verbatim}
\begin{table}[htbp]
\centering
\begin{tabular}{ll}
Command & Meaning \\
\hline
\LaTeX & Logo LaTeX \\
\TeX & Logo TeX
\end{tabular}
\caption{A small table}
\label{tab:small}
\end{table}
\end{verbatim}

Các tham số vị trí thường gặp:
\begin{center}
\begin{tabular}{>{\ttfamily}ll}
h & gần vị trí hiện tại;\\
t & đầu trang;\\
b & cuối trang;\\
p & trang chỉ chứa vật thể nổi;\\
! & nới lỏng một số ràng buộc.
\end{tabular}
\end{center}

Đừng đặt quá nhiều vật thể nổi liên tiếp mà không có văn bản xen giữa; \LaTeX\
cần đủ ngữ cảnh để bố trí chúng.

\subsection{Bảo vệ lệnh mong manh}

Một số lệnh không an toàn khi xuất hiện trong ``đối số chuyển động'', chẳng hạn
tiêu đề mục, chú thích hình hoặc mục lục. Nếu gặp lỗi khó hiểu, có thể cần
\cmd{protect}:
\begin{verbatim}
\section{Learning \protect\LaTeX}
\end{verbatim}

Trong \LaTeXe, nhiều lệnh đã được làm bền hơn, nhưng khái niệm lệnh mong manh
vẫn quan trọng khi làm việc với macro cũ.

\section{Công thức toán học}

\subsection{Kiến thức cơ bản}

\LaTeX\ có chế độ văn bản và chế độ toán. Công thức trong dòng đặt giữa
\verb|$...$| hoặc \verb|\(...\)|:
\begin{verbatim}
If $a^2+b^2=c^2$, the triangle is right-angled.
\end{verbatim}

Công thức hiển thị đặt trong \verb|\[...\]|:
\begin{verbatim}
\[
  a^2+b^2=c^2.
\]
\end{verbatim}

Không nên dùng \verb|$$...$$| trong \LaTeX\ hiện đại, vì nó bỏ qua một số cơ
chế khoảng cách và tương thích của \LaTeX.

Chỉ số trên và dưới dùng \verb|^| và \verb|_|:
\[
  x_i^2,\qquad a_{ij},\qquad e^{x+y}.
\]
Nếu chỉ số có nhiều hơn một ký tự, đặt trong ngoặc nhọn.

\subsection{Nhóm trong chế độ toán}

Dấu ngoặc nhọn nhóm các phần tử:
\begin{verbatim}
$x^2$ is different from $x^{2n}$.
\end{verbatim}

Phân số dùng \cmd{frac}:
\[
  \frac{x+y}{2},\qquad \frac{1}{1+\frac{1}{x}}.
\]

Căn dùng \cmd{sqrt}:
\[
  \sqrt{x},\qquad \sqrt[n]{x}.
\]

\subsection{Xây dựng khối công thức}

Các toán tử lớn:
\[
  \sum_{i=1}^{n} i = \frac{n(n+1)}{2},\qquad
  \int_{0}^{1} x^2\,dx = \frac{1}{3}.
\]

Dấu ngoặc tự co giãn dùng \cmd{left} và \cmd{right}:
\[
  \left(1+\frac{1}{1-x^2}\right)^3.
\]

Ma trận có thể viết bằng các môi trường của \pkg{amsmath}:
\[
\begin{pmatrix}
1 & 2\\
3 & 4
\end{pmatrix}
\qquad
\begin{bmatrix}
a & b\\
c & d
\end{bmatrix}.
\]

Hệ phương trình hoặc công thức nhiều dòng dùng \env{align}:
\begin{verbatim}
\begin{align}
f(x) &= x^2+2x+1\\
     &= (x+1)^2.
\end{align}
\end{verbatim}

Trong tài liệu này:
\begin{align}
f(x) &= x^2+2x+1\\
     &= (x+1)^2.
\end{align}

Nếu không muốn đánh số, dùng \env{align*}.

\subsection{Khoảng trắng trong toán học}

\TeX\ tự đặt phần lớn khoảng trắng toán học. Khi cần can thiệp, có các lệnh:
\begin{center}
\begin{tabular}{>{\ttfamily}ll}
\textbackslash , & khoảng nhỏ;\\
\textbackslash : & khoảng vừa;\\
\textbackslash ; & khoảng rộng;\\
\textbackslash ! & khoảng âm;\\
\textbackslash quad & khoảng một em;\\
\textbackslash qquad & khoảng hai em.
\end{tabular}
\end{center}

Ví dụ đúng hơn khi viết vi phân:
\[
  \int_0^1 x^2\,dx.
\]

Chữ trong công thức nên đặt bằng \cmd{text} hoặc \cmd{mathrm} tùy mục đích:
\[
  x_{\text{max}},\qquad \sin x,\qquad \mathrm{d}x.
\]

\subsection{Canh dọc}

\env{array} trong toán học tương tự \env{tabular}:
\[
\begin{array}{rcl}
x+y&=&7,\\
2x-y&=&4.
\end{array}
\]

Môi trường \env{cases} rất tiện cho hàm từng phần:
\[
|x| =
\begin{cases}
x, & x\ge 0,\\
-x, & x<0.
\end{cases}
\]

\subsection{Bóng ma}

Lệnh \cmd{phantom} tạo một hộp vô hình có kích thước như nội dung của nó. Nó
hữu ích khi muốn căn các biểu thức:
\[
\begin{array}{r}
\sqrt{\,x+1\,}\\
\sqrt{\phantom{x}+1}
\end{array}
\]

\cmd{hphantom} chỉ giữ chiều ngang, \cmd{vphantom} chỉ giữ chiều cao và chiều
sâu.

\subsection{Cỡ chữ toán học}

Trong công thức, \LaTeX\ tự chọn cỡ chữ phù hợp với bối cảnh: hiển thị, trong
dòng, chỉ số trên và chỉ số lồng sâu. Có thể ép bằng
\cmd{displaystyle}, \cmd{textstyle}, \cmd{scriptstyle} và
\cmd{scriptscriptstyle}, nhưng nên dùng ít.

Ví dụ:
\[
  \textstyle \sum_{i=1}^{n} i
  \qquad
  \displaystyle \sum_{i=1}^{n} i.
\]

\subsection{Định lý, định nghĩa}

Định lý và định nghĩa nên dùng môi trường được định nghĩa bằng \cmd{newtheorem}
để đánh số nhất quán.

\begin{verbatim}
\newtheorem{theorem}{Theorem}
\newtheorem{definition}{Definition}
\end{verbatim}

\begin{definition}
Một số nguyên \(p>1\) là số nguyên tố nếu nó chỉ có hai ước dương là \(1\) và
chính nó.
\end{definition}

\begin{theorem}
Có vô hạn số nguyên tố.
\end{theorem}

\subsection{Ký hiệu đậm}

Trong toán học, \cmd{textbf} không phải lúc nào cũng cho kết quả mong muốn.
Dùng \cmd{mathbf} cho chữ Latin đậm, hoặc \cmd{boldsymbol} nếu gói phù hợp có
sẵn. Với \pkg{unicode-math}, nhiều bảng chữ toán học có thể được chọn trực
tiếp:
\[
  \mathbf{x},\qquad \mathcal{F},\qquad \mathbb{R}.
\]

\subsection{Bảng ký hiệu toán học}

Các bảng dưới đây không cố liệt kê mọi ký hiệu của \LaTeX, nhưng giữ tinh thần
của bản gốc: cho người mới một nơi tra nhanh các lệnh thường dùng.

\subsubsection{Dấu phụ trong toán học}

\begin{center}
\begin{tabular}{llll}
\(\hat{a}\) & \cmd{hat} &
\(\bar{a}\) & \cmd{bar}\\
\(\tilde{a}\) & \cmd{tilde} &
\(\vec{a}\) & \cmd{vec}\\
\(\dot{a}\) & \cmd{dot} &
\(\ddot{a}\) & \cmd{ddot}\\
\(\widehat{AB}\) & \cmd{widehat} &
\(\overline{AB}\) & \cmd{overline}
\end{tabular}
\end{center}

\subsubsection{Chữ Hy Lạp}

\begin{center}
\begin{tabular}{llllllll}
\(\alpha\)&\cmd{alpha}&
\(\beta\)&\cmd{beta}&
\(\gamma\)&\cmd{gamma}&
\(\delta\)&\cmd{delta}\\
\(\epsilon\)&\cmd{epsilon}&
\(\varepsilon\)&\cmd{varepsilon}&
\(\zeta\)&\cmd{zeta}&
\(\eta\)&\cmd{eta}\\
\(\theta\)&\cmd{theta}&
\(\vartheta\)&\cmd{vartheta}&
\(\lambda\)&\cmd{lambda}&
\(\mu\)&\cmd{mu}\\
\(\pi\)&\cmd{pi}&
\(\varpi\)&\cmd{varpi}&
\(\rho\)&\cmd{rho}&
\(\sigma\)&\cmd{sigma}\\
\(\phi\)&\cmd{phi}&
\(\varphi\)&\cmd{varphi}&
\(\omega\)&\cmd{omega}&
\(\Omega\)&\cmd{Omega}
\end{tabular}
\end{center}

\subsubsection{Quan hệ và phép toán}

\begin{center}
\begin{tabular}{llllllll}
\(\le\)&\cmd{le}&
\(\ge\)&\cmd{ge}&
\(\ne\)&\cmd{ne}&
\(\approx\)&\cmd{approx}\\
\(\equiv\)&\cmd{equiv}&
\(\sim\)&\cmd{sim}&
\(\subset\)&\cmd{subset}&
\(\subseteq\)&\cmd{subseteq}\\
\(\in\)&\cmd{in}&
\(\notin\)&\cmd{notin}&
\(\cap\)&\cmd{cap}&
\(\cup\)&\cmd{cup}\\
\(\times\)&\cmd{times}&
\(\cdot\)&\cmd{cdot}&
\(\oplus\)&\cmd{oplus}&
\(\otimes\)&\cmd{otimes}
\end{tabular}
\end{center}

\subsubsection{Mũi tên, dấu ngoặc và ký hiệu khác}

\begin{center}
\begin{tabular}{llll}
\(\leftarrow\) & \cmd{leftarrow} &
\(\rightarrow\) & \cmd{rightarrow}\\
\(\leftrightarrow\) & \cmd{leftrightarrow} &
\(\Longrightarrow\) & \cmd{Longrightarrow}\\
\(\mapsto\) & \cmd{mapsto} &
\(\infty\) & \cmd{infty}\\
\(\partial\) & \cmd{partial} &
\(\nabla\) & \cmd{nabla}\\
\(\forall\) & \cmd{forall} &
\(\exists\) & \cmd{exists}
\end{tabular}
\end{center}

Dấu ngoặc có thể co giãn:
\[
\left\langle \frac{x}{y} \right\rangle,\qquad
\left\{ x\in\mathbb{R}\mid x^2<2 \right\}.
\]

\section{Chức năng đặc biệt}

\subsection{Chèn hình EPS}

Bản gốc trình bày cách chèn hình EPS bằng gói \pkg{graphicx}, vốn rất phổ biến
trong quy trình \LaTeX\ tạo DVI rồi chuyển sang PostScript hoặc PDF. Với quy
trình hiện đại tạo PDF trực tiếp, \pkg{graphicx} vẫn là gói chuẩn, nhưng định
dạng hình thường là PDF, PNG hoặc JPEG.

Ví dụ nguồn:
\begin{verbatim}
\usepackage{graphicx}

\begin{figure}[htbp]
  \centering
  \includegraphics[width=.8\linewidth]{figure-file}
  \caption{A sample figure}
  \label{fig:demo}
\end{figure}
\end{verbatim}

Trong bản dịch này không chèn hình thật để giữ tập tin tự chứa, nhưng các tùy
chọn quan trọng của \cmd{includegraphics} gồm:

\begin{center}
\begin{tabular}{>{\ttfamily}ll}
width & đặt chiều rộng;\\
height & đặt chiều cao;\\
scale & phóng to hoặc thu nhỏ theo tỷ lệ;\\
angle & xoay hình;\\
trim, clip & cắt mép hình.
\end{tabular}
\end{center}

\subsection{Tài liệu tham khảo}

Cách đơn giản nhất là dùng môi trường \env{thebibliography}:
\begin{verbatim}
\begin{thebibliography}{9}
\bibitem{lamport}
Leslie Lamport.
\newblock \emph{\LaTeX: A Document Preparation System}.
\newblock Addison-Wesley, 2nd edition, 1994.
\end{thebibliography}
\end{verbatim}

Trích dẫn bằng \cmd{cite}:
\begin{verbatim}
See Lamport~\cite{lamport}.
\end{verbatim}

Với tài liệu lớn, \BibTeX\ hoặc các hệ mới hơn như \texttt{biber} giúp quản lý
cơ sở dữ liệu tài liệu tham khảo riêng. Quy trình \BibTeX\ cổ điển là:
\begin{verbatim}
latex document
bibtex document
latex document
latex document
\end{verbatim}

\subsection{Chỉ mục}

Chỉ mục giúp người đọc tìm khái niệm trong sách. Quy trình cơ bản:
\begin{verbatim}
\usepackage{makeidx}
\makeindex

In the document body:
\index{LaTeX}

Where the index should be printed:
\printindex
\end{verbatim}

Sau khi biên dịch \LaTeX, chạy \MakeIndex\ để xử lý tập tin chỉ mục, rồi biên
dịch lại. Một mục chỉ mục có thể có khóa phụ:
\begin{verbatim}
\index{mathematics!formula}
\end{verbatim}

\subsection{Tùy chỉnh đầu trang và chân trang}

Gói \pkg{fancyhdr} thường được dùng để chỉnh đầu trang và chân trang:
\begin{verbatim}
\usepackage{fancyhdr}
\pagestyle{fancy}
\fancyhf{}
\fancyhead[LE,RO]{\thepage}
\fancyhead[LO]{\rightmark}
\fancyhead[RE]{\leftmark}
\end{verbatim}

Khi tùy chỉnh, hãy giữ trang dễ đọc. Đầu trang quá nhiều thông tin dễ làm phân
tán chú ý khỏi nội dung chính.

\subsection{Gói \pkg{verbatim}}

Môi trường \env{verbatim} chuẩn đủ cho nhiều trường hợp, nhưng gói
\pkg{verbatim} bổ sung tiện ích như chú thích cả một khối:
\begin{verbatim}
\begin{comment}
This whole block is ignored.
\end{comment}
\end{verbatim}

Nhiều tài liệu hiện đại dùng \pkg{listings} hoặc \pkg{minted} để tô sáng mã,
nhưng chúng không thuộc phạm vi của bản giới thiệu cổ điển này.

\subsection{Tải và cài đặt gói \LaTeX}

Các gói \LaTeX\ thường được phân phối qua \CTAN\ và đi kèm các bản phân phối
như TeX Live hoặc MiKTeX. Với TeX Live, trình quản lý \texttt{tlmgr} có thể cài
gói mới nếu hệ thống cho phép:
\begin{verbatim}
tlmgr install package-name
\end{verbatim}

Nếu cài thủ công, tập tin \file{.sty} phải nằm ở nơi \TeX\ tìm thấy, chẳng hạn
trong cây \file{texmf} cá nhân. Sau khi thêm tập tin vào cây hệ thống cũ, có
thể cần cập nhật cơ sở dữ liệu tên tập tin bằng \texttt{mktexlsr}. Cách chính
xác phụ thuộc bản phân phối và quyền trên máy.

\section{Tùy biến \LaTeX}

\subsection{Tạo lệnh, môi trường và gói mới}

\subsubsection{Tạo lệnh mới}

Lệnh mới được định nghĩa bằng \cmd{newcommand}:
\begin{verbatim}
\newcommand{\R}{\mathbb{R}}
\newcommand{\vect}[1]{\mathbf{#1}}
\newcommand{\norm}[1]{\left\lVert #1 \right\rVert}
\end{verbatim}

Tham số được đánh số \verb|#1|, \verb|#2|, v.v. Ví dụ:
\[
  \mathbb{R},\qquad \mathbf{x}.
\]

Nếu muốn định nghĩa lại lệnh đã tồn tại, dùng \cmd{renewcommand}. Hãy cẩn thận:
định nghĩa lại lệnh chuẩn có thể làm hỏng gói khác hoặc thay đổi kết quả ở
nhiều nơi không ngờ tới.

\subsubsection{Tạo môi trường mới}

Môi trường mới dùng \cmd{newenvironment}:
\begin{verbatim}
\newenvironment{important}
  {\begin{quote}\bfseries}
  {\end{quote}}
\end{verbatim}

Phần thứ hai là mã chạy lúc bắt đầu môi trường, phần thứ ba là mã chạy lúc kết
thúc môi trường.

\subsubsection{Tạo gói của riêng bạn}

Khi lời mở đầu của nhiều tài liệu lặp đi lặp lại, nên đưa các định nghĩa chung
vào một tập tin \file{.sty}. Ví dụ tối thiểu:
\begin{verbatim}
\NeedsTeXFormat{LaTeX2e}
\ProvidesPackage{mypackage}[2026/07/01 My macros]

\newcommand{\R}{\mathbb{R}}
\endinput
\end{verbatim}

Sau đó nạp bằng:
\begin{verbatim}
\usepackage{mypackage}
\end{verbatim}

\subsection{Phông chữ và cỡ chữ}

\subsubsection{Lệnh đổi kiểu chữ}

\LaTeX\ phân biệt lệnh khai báo và lệnh có đối số. Dạng có đối số thường an
toàn hơn cho thay đổi cục bộ:

\begin{center}
\begin{tabular}{lll}
\cmd{textrm} & \textrm{roman} & chữ có chân;\\
\cmd{textsf} & \textsf{sans serif} & chữ không chân;\\
\cmd{texttt} & \texttt{typewriter} & chữ đơn cách;\\
\cmd{textmd} & \textmd{medium} & nét thường;\\
\cmd{textbf} & \textbf{bold} & đậm;\\
\cmd{textup} & \textup{upright} & đứng;\\
\cmd{textit} & \textit{italic} & nghiêng;\\
\cmd{textsl} & \textsl{slanted} & xiên;\\
\cmd{textsc} & \textsc{Small Caps} & chữ hoa nhỏ.
\end{tabular}
\end{center}

Dạng khai báo như \cmd{itshape}, \cmd{bfseries}, \cmd{sffamily} có hiệu lực đến
hết nhóm hiện tại:
\begin{verbatim}
{\bfseries Only this part is bold.}
\end{verbatim}

Cỡ chữ chuẩn:
\begin{center}
\begin{tabular}{>{\ttfamily}ll}
\textbackslash tiny & rất nhỏ;\\
\textbackslash scriptsize & cỡ chỉ số;\\
\textbackslash footnotesize & cỡ chú thích;\\
\textbackslash small & nhỏ;\\
\textbackslash normalsize & bình thường;\\
\textbackslash large, \textbackslash Large, \textbackslash LARGE & lớn dần;\\
\textbackslash huge, \textbackslash Huge & rất lớn.
\end{tabular}
\end{center}

\subsubsection{Cảnh báo về lạm dụng phông chữ}

Bản gốc đùa bằng câu ``Danger, Will Robinson, Danger'' để nhắc rằng quá nhiều
kiểu chữ làm tài liệu kém chuyên nghiệp. Phông chữ nên phục vụ cấu trúc: tiêu
đề, thuật ngữ, nhấn mạnh, mã lệnh. Nếu mỗi đoạn lại đổi phông, người đọc sẽ
khó nhận ra đâu là thông tin quan trọng.

\subsubsection{Khuyến nghị}

Hãy dùng lệnh logic trước: \cmd{emph} cho nhấn mạnh, \cmd{section} cho tiêu đề,
môi trường định lý cho định lý. Chỉ dùng lệnh hình thức như \cmd{textbf} hoặc
\cmd{large} khi thật sự muốn kiểm soát ngoại hình cục bộ.

\subsection{Khoảng cách giữa các đối tượng}

\subsubsection{Khoảng cách dòng}

Khoảng cách dòng có thể đổi bằng \cmd{linespread} trong lời mở đầu:
\begin{verbatim}
\linespread{1.3}
\end{verbatim}
Tài liệu học thuật đôi khi yêu cầu giãn dòng, nhưng tăng quá nhiều sẽ làm trang
mất nhịp đọc.

\subsubsection{Định dạng đoạn}

Hai tham số thường gặp là \cmd{parindent} và \cmd{parskip}:
\begin{verbatim}
\setlength{\parindent}{1.5em}
\setlength{\parskip}{0.4em}
\end{verbatim}

Thông thường nên chọn một trong hai cách: thụt đầu dòng hoặc thêm khoảng cách
giữa đoạn. Dùng cả hai quá mạnh có thể làm văn bản rời rạc.

\subsubsection{Khoảng ngang}

Khoảng ngang thủ công:
\begin{verbatim}
\hspace{1cm}
\hspace*{1cm}
\end{verbatim}
Dạng có sao giữ khoảng trắng ngay cả ở đầu dòng, nơi khoảng trắng thường có thể
bị bỏ qua.

\cmd{stretch} tạo khoảng co giãn:
\begin{verbatim}
Left\hspace{\stretch{1}}Right
\end{verbatim}

\subsubsection{Khoảng dọc}

Khoảng dọc dùng \cmd{vspace}:
\begin{verbatim}
\vspace{1cm}
\vspace*{1cm}
\end{verbatim}

Trong văn bản bình thường, ít khi cần \cmd{vspace}. Nếu phải dùng nhiều, có lẽ
cấu trúc tài liệu hoặc lớp tài liệu chưa phù hợp.

\subsection{Bố cục trang}

Trang \LaTeX\ có nhiều tham số: chiều rộng vùng chữ, chiều cao vùng chữ, lề
trái, lề trên, khoảng đầu trang, khoảng chân trang. Bản gốc minh họa các tham
số như \cmd{textwidth}, \cmd{textheight}, \cmd{oddsidemargin},
\cmd{evensidemargin}, \cmd{topmargin}, \cmd{headheight}, \cmd{headsep} và
\cmd{footskip}.

Trong tài liệu hiện đại, gói \pkg{geometry} là cách thuận tiện hơn để đặt lề:
\begin{verbatim}
\usepackage[a4paper,margin=2.6cm]{geometry}
\end{verbatim}

Nếu tự đặt tham số, nhớ rằng một số độ dài của \LaTeX\ được tính từ mốc nội bộ,
không phải trực tiếp từ mép giấy. Vì vậy chỉnh từng tham số riêng lẻ dễ gây
nhầm lẫn.

\subsection{Chi tiết hơn về độ dài}

Độ dài trong \LaTeX\ có đơn vị. Một số đơn vị thường dùng:

\begin{center}
\begin{tabular}{>{\ttfamily}ll}
pt & point, đơn vị in ấn;\\
mm, cm & milimét, xentimét;\\
in & inch;\\
em & xấp xỉ chiều rộng chữ M của phông hiện tại;\\
ex & xấp xỉ chiều cao chữ x;\\
mu & đơn vị nhỏ dùng trong toán học.
\end{tabular}
\end{center}

Đặt độ dài:
\begin{verbatim}
\setlength{\textwidth}{15cm}
\addtolength{\parskip}{0.2em}
\end{verbatim}

Một độ dài có thể lưu trong thanh ghi riêng:
\begin{verbatim}
\newlength{\mylength}
\setlength{\mylength}{2cm}
\end{verbatim}

\subsection{Hộp}

\TeX\ xây trang từ các hộp. Người dùng thường gặp các lệnh:
\begin{verbatim}
\mbox{no line break}
\makebox[3cm][c]{centered}
\framebox[3cm][c]{framed}
\parbox{5cm}{A paragraph in a box of fixed width.}
\end{verbatim}

\cmd{makebox} nhận chiều rộng và vị trí tùy chọn: \texttt{l} căn trái,
\texttt{c} căn giữa, \texttt{r} căn phải, \texttt{s} giãn đều. \cmd{parbox}
cho phép văn bản trong hộp tự ngắt dòng theo chiều rộng đã cho.

\subsection{Thước và thanh chống}

\cmd{rule} tạo một hình chữ nhật đen:
\begin{verbatim}
\rule{3cm}{0.4pt}
\end{verbatim}
Kết quả là \rule{3cm}{0.4pt}.

Một ``thanh chống'' là một hộp vô hình dùng để ép chiều cao hoặc chiều sâu:
\begin{verbatim}
\rule{0pt}{2.4ex}
\end{verbatim}
Kỹ thuật này hữu ích trong bảng khi muốn các hàng cao đều hơn mà không vẽ gì
thêm.

\section{Phụ lục thực hành}

Phần này gom các ghi chú thực hành mà người đọc thường cần sau khi đã đọc năm
chương chính. Nó không thay thế chỉ mục đầy đủ của bản gốc, nhưng bổ sung một
lớp tra cứu nhanh cho quy trình soạn tài liệu kỹ thuật bằng \LaTeX\ hiện đại,
đặc biệt khi viết tiếng Việt và biên dịch bằng LuaLaTeX.

\subsection{Quy trình biên dịch hiện đại}

Với tài liệu đơn giản, một lần chạy LuaLaTeX có thể đủ:
\begin{verbatim}
lualatex document.tex
\end{verbatim}
Nhưng tài liệu có mục lục, tham chiếu chéo, danh sách hình bảng, thư mục hoặc
chỉ mục thường cần nhiều bước. Một chu trình đầy đủ hơn là:
\begin{enumerate}
  \item chạy \texttt{lualatex} để tạo \file{.aux}, \file{.toc}, \file{.idx} và
        các tập tin phụ;
  \item nếu dùng \BibTeX, chạy \texttt{bibtex document};
  \item nếu dùng chỉ mục, chạy \texttt{makeindex document};
  \item chạy \texttt{lualatex} thêm một hoặc hai lần để nhãn, số trang, mục lục
        và liên kết ổn định.
\end{enumerate}

Các công cụ như \texttt{latexmk} tự động hóa chu trình này:
\begin{verbatim}
latexmk -lualatex document.tex
\end{verbatim}
Khi biên dịch trong một dự án lớn, nên để tập tin sinh ra nằm trong thư mục
build riêng, như cách thư viện này sinh PDF vào \file{vi/build/...}. Điều đó
giữ nguồn \file{.tex} sạch và giúp dễ xóa toàn bộ sản phẩm tạm khi cần.

\subsection{Đọc lỗi biên dịch}

Thông báo lỗi của \TeX\ có thể dài, nhưng thường có cấu trúc lặp lại. Dòng bắt
đầu bằng dấu chấm than là lỗi chính; số dòng gần đó chỉ vị trí \TeX\ phát hiện
vấn đề, không phải lúc nào cũng là nguyên nhân thật.

\begin{center}
\begin{tabular}{>{\ttfamily}p{0.35\linewidth}p{0.52\linewidth}}
Undefined control sequence & Sai tên lệnh hoặc thiếu gói định nghĩa lệnh đó.\\
Missing \$ inserted & Đưa ký hiệu toán vào văn bản thường, hoặc quên đóng/mở
  chế độ toán.\\
Extra \} & Thừa ngoặc nhọn đóng hoặc thiếu ngoặc mở trước đó.\\
Runaway argument & Một tham số bắt buộc chưa được đóng, thường do thiếu
  \texttt{\}} hoặc môi trường chưa kết thúc.\\
Environment ... undefined & Sai tên môi trường hoặc thiếu gói cung cấp môi
  trường.\\
File ... not found & Đường dẫn hình/gói/tập tin nhập không đúng hoặc chưa cài
  gói.
\end{tabular}
\end{center}

Khi lỗi khó hiểu, hãy thu nhỏ ví dụ. Sao chép phần lời mở đầu cần thiết và một
đoạn văn tối thiểu vào tập tin mới; nếu lỗi biến mất, thêm dần các phần cho đến
khi tìm được lệnh gây lỗi. Đây là cách làm nhanh hơn nhiều so với đọc toàn bộ
nhật ký biên dịch của một tài liệu dài.

\subsection{Tiếng Việt và Unicode}

Với tài liệu tiếng Việt hiện đại, LuaLaTeX hoặc XeLaTeX thường thuận tiện hơn
quy trình \texttt{latex} cổ điển vì chúng đọc trực tiếp Unicode và dùng phông
OpenType qua \pkg{fontspec}. Một lời mở đầu tối thiểu có thể là:
\begin{verbatim}
\usepackage{fontspec}
\setmainfont{TeX Gyre Termes}
\setmonofont{DejaVu Sans Mono}[Scale=MatchLowercase]
\end{verbatim}

Các điểm cần chú ý:
\begin{itemize}
  \item lưu nguồn bằng UTF-8, không trộn nhiều bảng mã trong cùng dự án;
  \item chọn phông có đủ dấu tiếng Việt và ký tự toán/ký hiệu cần dùng;
  \item với toán học, dùng \pkg{unicode-math} nếu muốn phông toán Unicode;
  \item trong \env{verbatim} hoặc mã nguồn, phông mono phải có đủ ký tự cần
        hiển thị;
  \item không dùng \pkg{inputenc} như giải pháp chính trong LuaLaTeX, vì nó
        thuộc quy trình \LaTeX\ truyền thống hơn.
\end{itemize}

Nếu một chữ có dấu biến thành ô vuông trong PDF, nguyên nhân thường là phông
không có glyph đó, không phải do \LaTeX\ không hiểu tiếng Việt. Nếu chữ hiện
sai ngay trong nguồn, hãy kiểm tra mã hóa tập tin trước khi sửa lệnh \LaTeX.

\subsection{Hình, bảng và vật thể nổi}

Hình và bảng nên được đặt trong môi trường nổi khi chúng có chú thích và được
tham chiếu từ văn bản:
\begin{verbatim}
\begin{figure}
  \centering
  \includegraphics[width=.8\linewidth]{figures/example.pdf}
  \caption{Một hình minh họa.}
  \label{fig:example}
\end{figure}
\end{verbatim}

Lệnh \cmd{label} nên đặt sau \cmd{caption} để số hình/bảng được ghi đúng. Với
bảng, nên ưu tiên cấu trúc rõ ràng, ít đường kẻ dọc, và canh cột theo nghĩa dữ
liệu. Gói \pkg{booktabs} thường cho bảng đẹp hơn:
\begin{verbatim}
\begin{tabular}{ll}
\toprule
Tên & Ý nghĩa\\
\midrule
n & Số phần tử\\
\bottomrule
\end{tabular}
\end{verbatim}

Nếu hình gốc là sơ đồ thuật toán đơn giản, có thể vẽ lại bằng TikZ để giữ PDF
sắc nét và văn bản tìm kiếm được. Nếu hình là ảnh chụp hoặc biểu đồ phức tạp,
chèn ảnh raster/vector gốc thường tốt hơn. Không nên để lại một chú thích hình
đứng một mình khi hình đã bị bỏ; hoặc khôi phục hình, hoặc chuyển nội dung hình
thành prose rõ ràng.

\subsection{Thư mục tài liệu tham khảo}

Với tài liệu nhỏ, môi trường \env{thebibliography} là đủ:
\begin{verbatim}
\begin{thebibliography}{9}
\bibitem{lamport}
Leslie Lamport. \emph{LaTeX: A Document Preparation System}.
\end{thebibliography}
\end{verbatim}
Trích dẫn bằng \cmd{cite}:
\begin{verbatim}
See Lamport~\cite{lamport}.
\end{verbatim}

Với tài liệu lớn hoặc nhiều nguồn, nên dùng \BibTeX\ hoặc \texttt{biblatex}.
Cơ sở dữ liệu \file{.bib} giúp tách dữ liệu nguồn khỏi kiểu trình bày. Dù dùng
công cụ nào, khóa trích dẫn nên ổn định và có ý nghĩa, ví dụ
\texttt{knuth1984texbook}, thay vì \texttt{ref1}.

\subsection{Chỉ mục}

Chỉ mục giúp người đọc tìm thuật ngữ trong tài liệu dài. Quy trình cơ bản:
\begin{verbatim}
\usepackage{makeidx}
\makeindex
...
term\index{term}
...
\printindex
\end{verbatim}
Sau lần chạy \LaTeX, dùng:
\begin{verbatim}
makeindex document
\end{verbatim}
rồi biên dịch lại để đưa chỉ mục đã sắp xếp vào PDF.

Khi làm chỉ mục tiếng Việt, cần thống nhất cách ghi mục từ: có dấu hay không
dấu, tiếng Việt hay thuật ngữ gốc, viết hoa hay viết thường. Một tài liệu kỹ
thuật thường dùng dạng ``thuật ngữ tiếng Việt (term tiếng Anh)'' ở lần xuất
hiện đầu, rồi chỉ mục cả hai nếu người đọc có thể tra theo một trong hai cách.

\subsection{Chuyển một tài liệu dài sang \LaTeX\ sạch}

Khi dịch hoặc phục dựng một PDF dài, không nên bắt đầu bằng việc gõ lại từng
trang theo hình thức. Cách bền hơn là dựng trước cấu trúc logic:
\begin{enumerate}
  \item xác định lớp tài liệu, phông, kích thước trang và quy ước tiêu đề;
  \item tạo cây thư mục cho nguồn, hình, bảng dữ liệu và sản phẩm build;
  \item chia tài liệu thành các tập tin nhỏ theo chương hoặc phần bằng
        \cmd{input} hoặc \cmd{include};
  \item đặt quy ước nhãn như \texttt{sec:}, \texttt{fig:}, \texttt{tab:},
        \texttt{eq:}, \texttt{thm:};
  \item biên dịch sớm một bộ khung nhỏ trước khi nhập toàn bộ nội dung.
\end{enumerate}

Một lời mở đầu nên chứa các quyết định toàn cục, không chứa chỉnh sửa cục bộ
cho từng trang. Nếu một kiểu hộp, định lý, bảng hoặc chú thích lặp lại nhiều
lần, hãy định nghĩa lệnh hoặc môi trường riêng. Ngược lại, nếu một chỉnh sửa
chỉ dùng một lần, cân nhắc viết lại nội dung bằng văn bản thay vì thêm macro
khó hiểu.

\subsection{Quy tắc xử lý hình khi dịch PDF}

Hình trong PDF gốc có ba loại. Mỗi loại nên xử lý khác nhau:
\begin{itemize}
  \item \term{sơ đồ đơn giản}: vẽ lại bằng TikZ hoặc bảng/mũi tên để chữ tìm
        kiếm được và phông khớp với bản dịch;
  \item \term{biểu đồ hoặc ảnh phức tạp}: trích xuất ảnh/vector gốc, đặt trong
        thư mục \file{figures/}, rồi dịch caption và phần giải thích xung quanh;
  \item \term{hình chỉ là trang trí}: bỏ nếu không mang thông tin kỹ thuật, nhưng
        không để lại caption mồ côi.
\end{itemize}

Nếu chưa thể phục dựng hình ngay, nên thay bằng một đoạn mô tả có nội dung:
hình minh họa quan hệ gì, trục hoặc ký hiệu nào quan trọng, và người đọc cần
rút ra kết luận gì. Không nên để một khung trống với dòng ``hình sẽ bổ sung'',
vì PDF biên dịch được nhưng chất lượng đọc thấp.

\subsection{Tham chiếu chéo trong bản dịch}

Khi dịch tài liệu dài, số chương, hình, bảng và công thức có thể thay đổi so
với bản gốc. Vì vậy nên dùng \cmd{label} và \cmd{ref} thay vì gõ số tay:
\begin{verbatim}
\section{Đường đi ngắn nhất}\label{sec:shortest-path}
...
Như đã thấy ở Mục~\ref{sec:shortest-path}, ...
\end{verbatim}

Nhãn nên ổn định theo ý nghĩa, không theo số thứ tự hiện tại. Nếu sau này thêm
một hình trước đó, mọi tham chiếu vẫn đúng. Với công thức quan trọng, dùng
\env{equation} hoặc \env{align} có nhãn; với công thức phụ, dùng dạng không
đánh số để giảm nhiễu.

\subsection{Kiểm soát chất lượng trước khi commit}

Một bản dịch \LaTeX\ nên vượt qua các kiểm tra tối thiểu:
\begin{itemize}
  \item biên dịch sạch bằng engine đã chọn, ở đây thường là LuaLaTeX;
  \item không còn lỗi ``undefined reference'' hoặc ``citation undefined'';
  \item PDF có số trang hợp lý và mở được bằng \texttt{pdfinfo};
  \item trích xuất văn bản bằng \texttt{pdftotext} không lộ ghi chú việc còn
        làm, chỗ giữ tạm, caption mồ côi hoặc đoạn tiếng Anh ngoài ví dụ/code
        có chủ ý;
  \item hình và bảng không tràn trang, không che văn bản, không bị thiếu file;
  \item nguồn không có khoảng trắng lỗi hoặc file build tạm bị commit nhầm.
\end{itemize}

Trong một repo nhiều tài liệu, nên kiểm tra từng PDF sau khi sửa một file và
thỉnh thoảng chạy kiểm tra toàn bộ thư mục build. Điều này bắt lỗi phụ thuộc
gói, phông, đường dẫn hình và các vấn đề chỉ xuất hiện khi build từ thư mục gốc.

\subsection{Bảng thuật ngữ rút gọn}

\begin{longtable}{p{0.28\linewidth}p{0.58\linewidth}}
\textbf{Thuật ngữ} & \textbf{Ghi chú}\\
\hline
lời mở đầu & Phần trước \cmd{begin\{document\}}, nơi chọn lớp, gói và thiết lập
toàn cục.\\
thân tài liệu & Phần giữa \cmd{begin\{document\}} và \cmd{end\{document\}}.\\
lớp tài liệu & Tập \file{.cls}, quyết định cấu trúc và bố cục lớn.\\
gói & Tập \file{.sty}, bổ sung lệnh hoặc môi trường.\\
môi trường & Cặp \cmd{begin\{...\}} và \cmd{end\{...\}} tạo một vùng có quy tắc
dàn trang riêng.\\
vật thể nổi & Hình hoặc bảng được \LaTeX\ tự đặt ở vị trí phù hợp.\\
nhãn & Khóa dùng với \cmd{label} và \cmd{ref} để tham chiếu chéo.\\
nhật ký & Tập \file{.log}, chứa cảnh báo, lỗi và thông tin biên dịch.\\
overfull box & Dòng hoặc hộp vượt quá chiều rộng cho phép.\\
underfull box & Dòng hoặc hộp bị giãn quá mức để lấp đầy không gian.\\
\end{longtable}

\section*{Kết luận}
\addcontentsline{toc}{section}{Kết luận}

\LaTeX\ thưởng cho người viết biết kiên nhẫn với cấu trúc. Ban đầu, các lệnh
và bước biên dịch có vẻ xa lạ. Nhưng khi tài liệu có nhiều công thức, mục lục,
tham chiếu, bảng, hình, chỉ mục hoặc nhiều chương, lợi ích của việc mô tả nội
dung bằng nguồn văn bản rõ ràng trở nên rất lớn.

Nguyên tắc thực dụng là: dùng mặc định của \LaTeX\ cho đến khi có lý do thật
sự để đổi; dùng lệnh logic thay vì chỉnh ngoại hình cục bộ; chia tài liệu lớn
thành phần nhỏ; và đọc nhật ký biên dịch khi có lỗi. Với những thói quen đó,
\LaTeX\ trở thành một công cụ ổn định để viết tài liệu kỹ thuật lâu dài.

\begin{thebibliography}{9}
\bibitem{lamport}
Leslie Lamport.
\newblock \emph{\LaTeX: A Document Preparation System}.
\newblock Addison-Wesley, second edition, 1994.

\bibitem{knuth}
Donald E. Knuth.
\newblock \emph{The \TeX book}.
\newblock Addison-Wesley, 1984.

\bibitem{companion}
Michel Goossens, Frank Mittelbach and Alexander Samarin.
\newblock \emph{The \LaTeX\ Companion}.
\newblock Addison-Wesley, 1994.

\bibitem{ctan}
Comprehensive \TeX\ Archive Network.
\newblock \url{https://www.ctan.org}.
\end{thebibliography}

\end{document}
